%%% FIELD proof-conditional-completeness-reduction
Because \(\mathcal O\) is compatible, \(\mathcal C(\mathcal O)\neq\varnothing\).  Fix an arbitrary reported endpoint \(e\in\mathcal E(\mathcal O)\).  By \(\mathrm{def:monotone\mbox{-}closure}\), the mark assigned by \(\operatorname{Cl}_{\mathcal R}(\mathcal O)\) at \(e\) is either a member of \(\{\mathtt{tail},\mathtt{arrow}\}\) or \(\mathtt{circle}\).  We treat these two exhaustive cases.

Suppose first that
\[
 \operatorname{Cl}_{\mathcal R}(\mathcal O)(e)=m,
 \qquad m\in\{\mathtt{tail},\mathtt{arrow}\}.
 \tag{1}
\]
The soundness condition in the theorem gives
\[
 \mu_\xi(e)=m
 \qquad\text{for every }\xi\in\mathcal C(\mathcal O).
 \tag{2}
\]
By the invariant-mark clause of \(\mathrm{def:canonical\mbox{-}si\mbox{-}pag}\), equation \((2)\) implies
\[
 \mathcal P_{\mathrm{SI}}(\mathcal O)(e)=m
 =\operatorname{Cl}_{\mathcal R}(\mathcal O)(e).
 \tag{3}
\]

Suppose instead that
\[
 \operatorname{Cl}_{\mathcal R}(\mathcal O)(e)=\mathtt{circle}.
 \tag{4}
\]
Apply the opposing-extension condition first with \(m=\mathtt{tail}\) and then with \(m=\mathtt{arrow}\).  For each \(m\in\{\mathtt{tail},\mathtt{arrow}\}\), it supplies a base model \(\xi_{e,m}\in\mathcal C(\mathcal O)\) and a graph \(B_{e,m}\) admissible for \(\mathrm{lem:graphical\mbox{-}replacement\mbox{-}preservation}\) such that
\[
 \mu_{\widetilde\xi}(e)=m
 \quad\text{for every }
 \widetilde\xi\in\operatorname{Real}(\xi_{e,m},B_{e,m}).
 \tag{5}
\]
The faithful-realization condition gives
\[
 \operatorname{Real}(\xi_{e,m},B_{e,m})\neq\varnothing
 \qquad\text{for each }m\in\{\mathtt{tail},\mathtt{arrow}\}.
 \tag{6}
\]
Choose
\[
 \widetilde\xi_m\in
 \operatorname{Real}(\xi_{e,m},B_{e,m})
 \qquad\text{for each }m\in\{\mathtt{tail},\mathtt{arrow}\}.
 \tag{7}
\]
Since each \(B_{e,m}\) is admissible, \(\mathrm{lem:replacement\mbox{-}preservation}\), whose graphical premise is supplied by \(\mathrm{lem:graphical\mbox{-}replacement\mbox{-}preservation}\), yields
\[
 \mathcal O(\widetilde\xi_m)=\mathcal O,
 \qquad
 \widetilde\xi_m\in\mathcal C_A(\mathcal O,\xi_{e,m})
 \subseteq\mathcal C(\mathcal O).
 \tag{8}
\]
Combining \((5)\)--\((8)\) gives two members of the same compatibility fiber with opposing marks:
\[
 \mu_{\widetilde\xi_{\mathtt{tail}}}(e)=\mathtt{tail},
 \qquad
 \mu_{\widetilde\xi_{\mathtt{arrow}}}(e)=\mathtt{arrow}.
 \tag{9}
\]
Thus the endpoint mark is not invariant over \(\mathcal C(\mathcal O)\).  The circle clause of \(\mathrm{def:canonical\mbox{-}si\mbox{-}pag}\) therefore gives
\[
 \mathcal P_{\mathrm{SI}}(\mathcal O)(e)=\mathtt{circle}
 =\operatorname{Cl}_{\mathcal R}(\mathcal O)(e).
 \tag{10}
\]

Equations \((3)\) and \((10)\) establish equality at every \(e\in\mathcal E(\mathcal O)\).  Both objects are partial mixed-graph families on the reported domain fixed by \(\mathcal O\), by \(\mathrm{def:canonical\mbox{-}si\mbox{-}pag}\) and \(\mathrm{def:monotone\mbox{-}closure}\).  Hence endpointwise equality implies
\[
 \operatorname{Cl}_{\mathcal R}(\mathcal O)
 =\mathcal P_{\mathrm{SI}}(\mathcal O).
 \tag{11}
\]

Finally, suppose that every rule premise and path predicate has the stated uniform polynomial-time decision procedure.  The list \(\mathcal R_{\mathrm{SI}}\) is finite and its schemata are independent of \(D\) and \(K\) by hypothesis.  Membership in \(\Lambda_{\mathrm{SI}}\), together with \(\mathrm{def:monotone\mbox{-}closure}\), ensures that each strict conclusion replaces at least one reported circle by a definite endpoint and that no definite endpoint is subsequently changed.  Therefore all hypotheses of \(\mathrm{prop:exact\mbox{-}conditional\mbox{-}polynomial\mbox{-}saturation}\) hold.  That proposition yields polynomial-time saturation and at most
\[
 (K+1)D(D-1)
\]
strict reported-endpoint updates.  The argument uses the faithful-realization and opposing-extension conditions only as explicit antecedents; it neither establishes them nor constructs the candidate rule list.  This proves the conditional claim.

